\section{The dyadic distribution estimate}\label{sec:dyadic}

The sparse estimate supplies a gain when every packet meets little measure.
For a general array, we first remove sparse packets and narrow angular
contributions.  The remaining packets lie in a controlled collection of
plates.  The main analytic step then replaces all intermediate packet scales
by a single family at a shorter scale, to which the bilinear frame estimate
applies.

\begin{theorem}[Dyadic distribution bound]\label{thm:distribution}
There are fixed $p>2$ and $\sigma>0$ with the following property.
Let $U$ be a finite sum of the atoms in \eqref{eq:src5}, of packet size
$L\geq1$, with $\sum_\alpha|c_\alpha|^2\leq1$.  Suppose that $\mu$ is
supported in a spatial box of side $O(L^2)$ and in $|t|\leq L^2$, that
$\mu(B_r)\leq mr^2$ for every ball of radius $r\geq1$, and that
\eqref{eq:src6} holds for every label, with $1\leq D\leq L$.
Then, for every $g>0$,
\begin{equation}\label{eq:src12}
 \mu(A)\leq C mL^4D^{-\sigma}g^{-p},
 \qquad A=\{z:|U(z)|\geq\lambda\},\qquad
 \lambda=gL^{-4/3}.
\end{equation}
The constant is uniform in $L,D,g,m$ and the size of the array.  It may
depend on the fixed packet class, including sufficiently many finite
profile derivative bounds.  For an arbitrary coefficient array, the right
side is multiplied by $\|c\|_{\ell^2}^p$.
\end{theorem}

We fix the packet class throughout the proof.  In particular, its frequency
square, phase-space counting constant, profile support, and required
derivative bounds are fixed.  Recursive cap scalings will be performed in
the original, unrotated coordinates, using squares and retaining only the
specified subarray.  This preserves the class, including its derivative
bounds.  Rotations used later are only in nonrecursive estimates.  The
regularity order can be chosen after all small exponents: the saving $c_s$
in Proposition~\ref{prop:sparse}, the geometric powers in
Proposition~\ref{prop:rich}, and the separation power in
Theorem~\ref{thm:frame} do not increase when additional profile regularity
is required.

\subsection{Parameters and recursive estimates}

Write $c_s$ for the saving in \eqref{eq:src7}, and $p_0>2$ for the exponent
in \eqref{eq:src9} and \eqref{eq:src10}. Choose
$2<p_{\max}<\min(p_0,3)$ and then $\eta>0$ so small that
\begin{equation}\label{eq:dyadic-gamma}
 \gamma_0=1-p_{\max}/3-p_{\max}\eta>0.
\end{equation}
We introduce the parameters used in the construction and work conditionally
on the inequalities below. Their compatible order of choice will be checked
after the estimates, when the losses determining them have been identified.
For fixed constants $B,C_a$ and a small exponent $u>0$, set
\begin{equation}\label{eq:dyadic-parameters}
 X=BD,\qquad D_*=BD^2,\qquad \delta_0\asymp X^{-u},
 \qquad a\asymp X^{C_a},\qquad h=L/a.
\end{equation}
Round $\delta_0$ and $h$ to convenient dyadic scales, changing only fixed
constants. The parameter $D_*$ will be used in the sparse-packet deletion;
$\delta_0$ is the upper scale for cap recursion. The shorter packet scale
$h$ will be used after the geometric reduction.

We split the proof at $D=L^\kappa$, with $\kappa>0$ sufficiently small.
In the range $D<L^\kappa$, we require $h$ to dominate the fixed powers
of $X$ occurring below, after enlarging a fixed base scale $L_0$. We also require
$2<p\leq p_{\max}$ and $0<\sigma\leq1$ with
\begin{equation}\label{eq:dyadic-exponent-conditions}
 \sigma+\frac{4(p-2)}{3\kappa}<c_s,\qquad
 C_fu+C_g(p-2)+\sigma<c_s,\qquad
 \sigma<u\gamma_0.
\end{equation}
Here $C_g$ is the fixed exponent separating the large levels
$g\geq X^{C_g}$ from the remaining levels, and $C_f$ is a fixed loss exponent, independent of $u$. The frame
application below will determine a sufficient value of $C_f$ for
\eqref{eq:dyadic-final-energy}. All three
inequalities have strict margins. We take $B$ sufficiently large for the
recursive terms to contract, and choose the overall induction constant
only after the nonrecursive estimates have been obtained.

Two elementary estimates start the induction.  At a point, compact support
allows only $O(L^2)$ labels, each atom is $O(L^{-1})$, and hence
\begin{equation}\label{eq:dyadic-sup}
 \|U\|_\infty\lesssim\|c\|_{\ell^2}.
\end{equation}
Also $\mu(\supp\mu)\lesssim mL^4$.  Thus all bounded scales are covered
by a constant depending on $L_0$.  When $D\geq L^\kappa$, unit-cube mass
is $O(m)$, so Proposition~\ref{prop:sparse} gives
\[
 \int|U|^2\,d\mu
 \leq C m^{1/3}\sum_q\mu(q)^{2/3}\sup_q|U|^2
 \lesssim_\kappa mL^{4/3}D^{-c_s}.
\]
Chebyshev therefore gives $CmL^4D^{-c_s}g^{-2}$.  A nonempty level set
has $g\lesssim L^{4/3}$ by \eqref{eq:dyadic-sup}; the first inequality
in \eqref{eq:dyadic-exponent-conditions} implies
$D^\sigma g^{p-2}\lesssim D^{c_s}$.  This proves
\eqref{eq:src12} in the sparse range independently of induction.

For $D<L^\kappa$, use outer induction on dyadic ranges of $L$ and inner
induction on the number of iterations of $d\mapsto Bd^2$ needed to reach
$L^\kappa$.  Thus the result is available at a fixed smaller packet
scale, and at the present scale with parameter $D_*$.  For sufficiently
large $L$ in this range, $D_*\leq L$; it suffices to have
$BL^{2\kappa}\leq L$.  Denote the inductive constant by $C_{\rm ind}$.
We will bound recursive contributions by a fixed fraction of
$C_{\rm ind}mL^4D^{-\sigma}g^{-p}$ and all other contributions by a
constant independent of $C_{\rm ind}$.

\begin{lemma}[Cap gain]\label{lem:dyadic-cap-gain}
For any subarray in a frequency square of side comparable to $r$, where
$L^{-1}\lesssim r\leq\delta_0$, outer induction gives
\begin{equation}\label{eq:dyadic-cap-gain}
 \mu\{|U_{\mathcal C}|\geq r^\eta\lambda\}
 \leq C_{\rm cap}C_{\rm ind}mL^4g^{-p}
 r^{1-p/3-p\eta}
 \Big(\sum_{\alpha\in\mathcal C}|c_\alpha|^2\Big)^{p/2}.
\end{equation}
The estimate holds for any submeasure of $\mu$, with the same upper ball
parameter.  At the smallest scales choose comparable squares so that
the child packet size is at least $1$.
\end{lemma}
\begin{proof}
For a square centered at $v$, use
$(x',t')=(r(x-2vt),r^2t)$, up to fixed normalization conventions.  The
child packet size is $L'=rL$, the field factor is $r$, and the pushed
measure has ball parameter $m'\leq Cm r^{-3}$.  The new threshold
parameter is comparable to $g r^{1/3+\eta}$, since
\[
 r^{-1}r^\eta gL^{-4/3}
   =g r^{1/3+\eta}(rL)^{-4/3}.
\]
After a fixed inflation of $m'$, covering tubes and their transverse
annuli supplies \eqref{eq:src6} with $D=1$; equivalently
$m'(L')^3\asymp mL^3$.  The time interval is still
$|t'|\leq(L')^2$.  Partition the transformed spatial range into boxes
of side comparable to $(L')^2$.  A compact child packet meets only
boundedly many such boxes during this time interval.  If their relevant
coefficient energies are $e_j$, then
$\sum_j e_j\lesssim\sum_{\mathcal C}|c_\alpha|^2$ and
$\sum_j e_j^{p/2}\leq(\sum_j e_j)^{p/2}$.  Consequently outer induction
costs no number of boxes.  Its dimensional factor is
\[
 (mr^{-3})(rL)^4(g r^{1/3+\eta})^{-p}
  =mL^4g^{-p}r^{1-p/3-p\eta},
\]
as asserted.  These parabolic transformations translate the pairs
$(b_\alpha/L,L\omega_\alpha)$ and leave the normalized profiles
unchanged.  Enclosing a specified cap in a comparable square does not
add any other labels.
\end{proof}

We need the same estimates when the output point determines a prefix of
the array.  The next elementary device avoids a loss depending on the
number of labels. It adapts the dyadic input-energy decomposition of
Christ and Kiselev \cite[Theorem~1.1 and its proof in Section~2]{CK01}
to the weak-type block estimates on restricted submeasures used here.
We give the proof for these hypotheses, rather than applying their
maximal theorem as a black box.

\begin{lemma}[Energy truncation]\label{lem:dyadic-prefix}
Let a finite coefficient array be ordered, with total energy $E$, and let
each output point select a prefix ending between whole coefficients.
Suppose the distribution estimate for every block, on the submeasure
where that block is selected, has right side
$K t^{-p}\|c_{\rm block}\|_{\ell^2}^p$, where $p>2$.
The field formed from the selected prefix then satisfies the same
estimate with right side $C_p Kt^{-p}E^{p/2}$.
\end{lemma}
\begin{proof}
Zero coefficients may be omitted, and $E=0$ is trivial.  Place successive
intervals of lengths $|c_\alpha|^2/E$ in $[0,1]$.  For a dyadic interval
$I\subset[0,1]$, apportion $c_\alpha$ by the fraction of its interval
lying in $I$.  If that fraction is $f_{\alpha,I}$, its apportioned
energy obeys
\[
 \sum_\alpha |f_{\alpha,I}c_\alpha|^2
 \leq\sum_\alpha f_{\alpha,I}|c_\alpha|^2=E|I|.
\]
These apportioned blocks satisfy the assumed distribution estimate, up
to a constant depending on $p$. Each is either a single apportioned
coefficient or the sum of a whole interior block and at most two
fractional endpoint coefficients. Homogeneity gives the estimate for
the fractional singletons. Split the threshold equally among these at
most three pieces and use
$\sum_i\|c_i\|_2^p\leq(\sum_i\|c_i\|_2^2)^{p/2}$;
the pieces have disjoint coefficient supports. The resulting bound uses
exactly the apportioned energy above.

Every prefix interval $[0,s)$ is a disjoint union of at most one dyadic
interval at each depth $k\geq0$.  Every coefficient contributing to a
used interval belongs to the prefix, because $s$ is between whole
coefficient intervals.  The sum of these block fields equals the prefix
field, also when the binary expansion is infinite; the original array
is finite and the coefficient fractions sum exactly.

Choose $0<\epsilon<(p/2-1)/p$ and positive thresholds
$t_k=c_\epsilon t2^{-\epsilon k}$ with $\sum_k t_k=t$.
If a prefix field exceeds $t$, some used block exceeds $t_k$.
There are at most $2^k$ blocks at depth $k$, each with energy at most
$E2^{-k}$.  Restricting the measure separately to points using a block,
and summing the distribution bounds, gives
\[
 3^pKt^{-p}E^{p/2}c_\epsilon^{-p}
 \sum_{k\geq0}2^{k(1-p/2+p\epsilon)}
 \leq C_pKt^{-p}E^{p/2}.
\]
This argument also applies inside each fixed frequency cap, with its
own total energy $E$.
\end{proof}

\subsection{Peeling to a common rich measure}

Starting with $A$, repeatedly perform the following two deletions.
First remove all labels whose smoothed mass against the measure on the
current set is less than $mL^3/D_*$.  Next remove every point at which
the remaining physically neighboring labels do not support a probability
law with direction-ball bounds
\begin{equation}\label{eq:dyadic-angular-law}
 \mathbb P\{\omega\in B_r(v)\}
 \leq C_{\rm ang}\delta_0^{-\eta}r^\eta,
 \qquad L^{-1}\leq r\lesssim1.
\end{equation}
Here a physical neighbor means that the point is within a fixed multiple
of $L$ of the packet trajectory; the multiple contains the compact
profile support.  Directions and frequency centers are interchangeable
up to fixed constants on the bounded frequency square.

The angular test can use a finite dyadic label tree, ending at scale
comparable to $L^{-1}$, so it depends measurably on one of finitely many
incidence patterns.  If a round deletes points but no subsequent label,
the angular test is unchanged and no further point can be deleted.
Since there are finitely many labels, the process terminates.  Write
$A_\infty$ for its stabilized set and $\mathcal R$ for its surviving
labels.  At termination all labels in $\mathcal R$ satisfy
\begin{equation}\label{eq:dyadic-rich-common}
 \int_{A_\infty}
 (1+|x-b_\alpha-2\omega_\alpha t|/L)^{-M}\,d\mu
 \geq mL^3/D_*;
\end{equation}
every point of this same set satisfies \eqref{eq:dyadic-angular-law}
using physically neighboring labels in $\mathcal R$.

To estimate what was deleted, order removed labels by their removal
round, breaking ties arbitrarily.  Each original point selects the
prefix removed before it exits, including all labels removed in its exit
round; a point of $A_\infty$ selects all removed labels.  Let $P(z)$ be
this prefix field.  For any binary block in
Lemma~\ref{lem:dyadic-prefix}, restrict the measure to points whose
prefix uses the block.  Such points were still present when each
contributing label was removed.  By monotonicity of the smoothed
integral, every label in that block satisfies \eqref{eq:src6} with
$D_*$.  Inner induction and Lemma~\ref{lem:dyadic-prefix} therefore
give, for any fixed small $c>0$,
\begin{equation}\label{eq:dyadic-prefix-exception}
 \mu\{|P|\geq c\lambda\}
 \leq C_p C_{\rm ind}mL^4D_*^{-\sigma}g^{-p}.
\end{equation}
This restriction of measure is essential: a removed label need not be
sparse on all of the original set.

For both the original field and this prefix field, also remove the
exceptional sets where a dyadic cap contribution at any scale
$L^{-1}\lesssim r\leq\delta_0$ exceeds $r^\eta\lambda$.
Lemma~\ref{lem:dyadic-cap-gain} and, for prefixes,
Lemma~\ref{lem:dyadic-prefix}, imply a total cost
\begin{equation}\label{eq:dyadic-cap-exceptions}
 C' C_{\rm ind}mL^4g^{-p}\delta_0^{\gamma_0}.
\end{equation}
Indeed cap energies at each scale partition the original energy, so
their $p/2$ powers sum to at most $1$, and
$1-p/3-p\eta\geq\gamma_0$ makes the dyadic scale sum geometric.

\begin{lemma}[The angular test cannot fail elsewhere]
\label{lem:dyadic-flow}
Outside the exceptions in \eqref{eq:dyadic-prefix-exception} and
\eqref{eq:dyadic-cap-exceptions}, an original point cannot exit in an
angular deletion.  On a surviving such point,
$|U_{\mathcal R}|\geq c\lambda$ for an absolute $c>0$.
\end{lemma}
\begin{proof}
At a proposed exit, subtract the selected prefix from $U$.  The
remaining field has modulus at least a fixed fraction of $\lambda$,
and every fine cap sum has modulus at most $2r^\eta\lambda$.
Only physically neighboring remaining labels can contribute at the
point.

Mark the terminal cells of the finite direction tree that contain such
a neighboring label.  Give unmarked terminal cells capacity zero,
fine cells of side $r\leq\delta_0$ capacity $r^\eta$, and coarser cells
capacity $1$.  A cut covering the marked leaves either contains a
coarse cell, and hence has capacity at least $1$, or consists of fine
cells.  In the latter case its cell fields sum to the remaining field,
so the triangle inequality gives total capacity at least an absolute
positive constant.  Nonneighboring labels in these cells contribute
zero at the point.

For completeness, the maximal mass that can pass through a node is
the minimum of its capacity and the sum of the maximal masses of its
children.  Induction up the finite tree shows that this is the minimum
cut capacity.  Thus a mass bounded below can be sent from the root to
marked leaves.  Choose a neighboring label in each such leaf and
normalize the resulting flow to a probability.  A fine cell of side $r$
then has probability $O(r^\eta)$.  A ball of radius $r\geq L^{-1}$ is
covered by boundedly many cells of comparable size.  At coarser scales
$r\geq\delta_0$, the desired bound follows from
$\delta_0^{-\eta}r^\eta\geq1$.  Taking $C_{\rm ang}$ large enough
gives \eqref{eq:dyadic-angular-law}, a contradiction to the exit test.
The assertion about the final field follows by subtracting its prefix.
\end{proof}

Relative to the target with constant $C_{\rm ind}$, the recursive costs
so far are bounded by fixed multiples of
\begin{equation}\label{eq:dyadic-contraction}
 (D_*/D)^{-\sigma}=X^{-\sigma},\qquad
 D^\sigma\delta_0^{\gamma_0}
 \lesssim D^\sigma X^{-u\gamma_0}.
\end{equation}
The later single-cap exceptions have the second form as well.  The last
inequality in \eqref{eq:dyadic-exponent-conditions} and a sufficiently
large $B$ make the sum of all these costs less than, say, half the
target.  Their constants include the fixed prefix constant $C_p$, but
do not include any of the final nonrecursive constants.

We apply geometry to $\mu|_{A_\infty}$ itself, before subtracting any of
the proof-only exceptional sets.  Equations
\eqref{eq:dyadic-angular-law} and \eqref{eq:dyadic-rich-common}
therefore hold for the same measure.  Subsequent upper bounds may use
the original $\mu$: its ball and sparsity upper bounds still hold for
every surviving label.

\subsection{Plates and coarse frequency intervals}

Divide physical coordinates by $L$, and divide $\mu|_{A_\infty}$ by
$mL^2$.  In the notation of Proposition~\ref{prop:rich}, $H=L$,
the richness is at least $L/D_*\geq L/X^2$, and the angular constant
in \eqref{eq:dyadic-angular-law} is at most a fixed power of $X$.
The required comparisons $X^C\ll H$ hold by the choice of $\kappa$.
Proposition~\ref{prop:rich} assigns every surviving tube to one of at
most $X^{C_1}$ plates of width $X^{C_1}$ in these coordinates, with
\[
 |n\cdot(2\omega,1)|\leq X^{C_1}/L.
\]
Enlarge the plates by a fixed polynomial factor to contain the physical
packet supports in the domain.  All assignments are disjoint at the
level of coefficient labels.

The spatial component $n_x$ of such a unit normal is bounded below.
Indeed $|\omega|\leq C$, so if $|n_x|$ were sufficiently small, then
$|n_t+2n_x\cdot\omega|$ would be bounded below by an absolute constant,
contrary to the displayed estimate for large $L$.  Thus the assigned
frequencies lie in a strip of width $X^{O(1)}/L$, which can be partitioned
into polynomially many strips of width $O(L^{-1})$.  If the total number
of resulting arrays is $N\leq X^{C_2}$, the triangle inequality and
Proposition~\ref{prop:strip} give
\begin{equation}\label{eq:dyadic-large-level}
 \mu\{|U_{\mathcal R}|\geq c\lambda\}
 \lesssim mL^4N^{p_0}g^{-p_0}
       \sum_{j=1}^N\|c_j\|_{\ell^2}^{p_0}
 \leq X^{C_3}mL^4g^{-p_0}.
\end{equation}
Choosing $C_g(p_0-p_{\max})>C_3+1$ makes this at most a fixed multiple
of the target whenever $g\geq X^{C_g}$.  For $g<D^{-1}$, total mass
suffices, since $D^{-\sigma}g^{-p}\geq D^{p-\sigma}\geq1$ after
decreasing $\sigma$ if necessary.  We henceforth work in
\begin{equation}\label{eq:dyadic-middle-levels}
 D^{-1}\leq g\leq X^{C_g}.
\end{equation}

Discard intersections of enlarged plates whose projective normal
directions differ by at least a constant multiple of $a/L$.  For two
plates of normalized width $w\leq X^{C_4}$ and normal separation
$\beta$, their intersection is in a box of length $O(L)$ and transverse
sides $O(w)$ and $O(w/\beta)$.  Its normalized mass is at most
$X^{O(1)}L/\beta$ by unit-ball covering.  Summing over the polynomial
plate list, the discarded mass in the original coordinates is at most
\begin{equation}\label{eq:dyadic-plate-discard}
 mL^4X^{C_5}/a.
\end{equation}
Choose $C_a>C_5+1+C_gp_{\max}$, with room for fixed constants.  Then
\eqref{eq:dyadic-plate-discard} is bounded by the target throughout
\eqref{eq:dyadic-middle-levels}.

Partition projective normal directions into cells of diameter comparable
to $a/L$, with bounded overlap on balls of that radius.  At a retained
point all active plate normals have pairwise distance $O(a/L)$.
Consequently only boundedly many normal cells contribute, and at least
one cluster field has modulus at least $c\lambda$.  Relative to a
representative normal in each nonempty cell, its label frequencies lie
in a strip of width $O(a/L)=O(h^{-1})$.  The representative still has
spatial component bounded below, and this follows by perturbing the
previous direction inequality by $O(a/L)$.  The clusters partition the
surviving coefficient energy.

For each cluster separately rotate the spatial coordinates and remove
its central normal frequency by a bounded shear and modulation.
Write the resulting coordinates as $(x,y,t)$, with $x$ tangential and
$y$ normal.  Partition tangential label frequencies into coarse
intervals $\tau$ of length comparable to $\delta_0$, then into fine
intervals $\theta$ of length comparable to $h^{-1}$, using nested grids
and consistent endpoint assignments.  Each coarse array fits in a
comparable two-dimensional cap, because $h^{-1}\ll\delta_0$.
In applying Lemma~\ref{lem:dyadic-cap-gain}, enclose that array in an
unrotated square in the original coordinates, keeping only its labels.
At each scale bounded overlap of the enclosing squares, together with
the partitioned coefficient energy across clusters, gives a total
exceptional cost
\[
 C''C_{\rm ind}mL^4g^{-p}\delta_0^{\gamma_0}
\]
for the sets on which a coarse cluster field exceeds
$\delta_0^\eta\lambda$.  This is the further recursive cost already
allowed in \eqref{eq:dyadic-contraction}.

At any remaining point where a cluster is high, at least
$c\delta_0^{-\eta}$ distinct coarse intervals satisfy
\begin{equation}\label{eq:dyadic-significant-intervals}
 |U_\tau|\geq c'\delta_0\lambda.
\end{equation}
To see this, there are $O(\delta_0^{-1})$ coarse intervals, so the sum
of those below the displayed small threshold has modulus less than
half the cluster threshold if $c'$ is sufficiently small.  Each of the
others has modulus at most $\delta_0^\eta\lambda$, forcing the stated
number.  Full partition intervals are used throughout; the complement
estimate never discards only part of an interval field.

\subsection{Flattening all intermediate scales}

The cluster frequencies now lie in a strip of width $O(h^{-1})$, whereas
Proposition~\ref{prop:strip} at packet size $L$ requires width $O(L^{-1})$.
Splitting into thinner strips produced the large-level estimate, but paid
a fixed power of $X$. To handle the remaining levels, we instead make one
smooth packet out of each entire fine frequency array on each short
physical cell. We will apply Theorem~\ref{thm:frame} directly to these
packets and sum their energies. The construction uses only the general
compact profiles in \eqref{eq:src5}; no equation or exact Fourier
support is imposed on those profiles.

For a fine interval $\theta$ centered at $v_\theta$, put
\begin{equation}\label{eq:dyadic-child-coordinates}
 z'_\theta=F_\theta(x,y,t)
   =\big((x-2v_\theta t)/h,\ y/h,\ t/h^2\big).
\end{equation}
The central normal frequency has already been removed.  The array is
contained in a cap of size $O(h^{-1})$ and therefore has the form
\begin{equation}\label{eq:dyadic-child-field}
 U_\theta(x,y,t)
 =h^{-1}e^{i(xv_\theta-tv_\theta^2)}U'_\theta(z'_\theta),
\end{equation}
up to a common harmless modulation.  Here $U'_\theta$ is an array of
packet size $a=L/h$, with bounded frequencies and the original
coefficient norm.

Choose a fixed smooth partition of unity $\{\psi_T\}$ on the unit
lattice in $\R^3$, with bounded overlap and fixed enlarged supports.
Take a finite derivative order $N$ exceeding all later Sobolev and frame
requirements, and let
\[
 s_{\theta,T}
   =h^{-1}\max_{|\beta|\leq N}
       \sup_{z'\in T^*}|\partial^\beta U'_\theta(z')|,
\]
where $T^*$ is a fixed sufficiently large enlargement of $T$.  Enlarging
the fixed constant in this definition would give the same estimates.
Fix also a large polynomial decay power $P$, larger than $M+3$ and all
tail moments that will occur, and define
\begin{equation}\label{eq:dyadic-density}
 d_\theta(T)=\frac{1}{mh^3}
   \int(1+\dist(F_\theta(z),T))^{-P}\,d\mu(z).
\end{equation}
We abbreviate this to $d(T)$ when the cluster and $\theta$ are specified.
All coordinate changes push forward the measure without changing its
mass.  Their ball constants are uniformly bounded after the displayed
factor $h^3$.  In particular $d(T)\lesssim1$ by annular covering.

Only cells whose partition supports meet the testing domain need be
retained.  The associated summand
\begin{equation}\label{eq:dyadic-composite-packet}
 e^{i(xv_\theta-tv_\theta^2)}
 h^{-1}\psi_T(z'_\theta)U'_\theta(z'_\theta)
\end{equation}
is a single composite packet: it has tangential width $O(h)$, time
length $O(h^2)$, localized dependence on $y/h$, and all required
normalized derivatives bounded by $Cs_{\theta,T}$.

There are two energy bounds to retain. On a time interval of length $h^2$,
spatial almost orthogonality controls the total squared amplitudes of the
composite packets. Over the full time range, the sparse estimate controls
the same amplitudes weighted by $d(T)^{2/3}$. The first bound will supply
the large threshold-to-envelope ratios required by the frame theorem;
the second will retain the factor $D^{-c_s}$ when the frame estimates are
summed.

\begin{lemma}[Flattened energy estimates]\label{lem:dyadic-flattening}
For a time interval $J$ of length $h^2$, let $\mathcal T(J)$ denote the
boundedly many native time-cell layers that can contribute on $J$.
Summing over the disjoint fine arrays, and also over clusters when
needed, one has
\begin{equation}\label{eq:src13}
 \begin{aligned}
 \sum_{\theta,\,T\in\mathcal T(J)}s_{\theta,T}^2
     &\lesssim h^{-2},\\
 \sum_{\theta,T}d(T)^{2/3}s_{\theta,T}^2
     &\lesssim_B h^{-2}a^{4/3}D^{-c_s}.
 \end{aligned}
\end{equation}
The right sides are multiplied by the total original coefficient
energy if it is not normalized to be at most $1$.
\end{lemma}
\begin{proof}
We first justify spatial almost orthogonality for arbitrary smooth
packet profiles.  At native time $t'=O(a^2)$, a child atom and all its
fixed derivatives have spatial support within $O(a)$ of its transported
intercept $b'_\alpha+2\omega'_\alpha t'$.  Spatial integration by
parts in the difference of the carrier frequencies gives Gram entries
bounded by
\begin{equation}\label{eq:dyadic-gram}
 C_N(1+a|\omega'_\alpha-\omega'_\beta|)^{-N},
\end{equation}
and the entries vanish unless the transported intercepts differ by
$O(a)$.  The joint centers
\[
 (b'_\alpha/a,a\omega'_\alpha)
 \longmapsto
 \big(b'_\alpha/a+2(t'/a^2)a\omega'_\alpha,
      a\omega'_\alpha\big)
\]
undergo a uniformly bounded shear.  Thus their phase-space counts remain
bounded.  Summing \eqref{eq:dyadic-gram} over these centers gives
uniformly summable Gram rows when $N$ is sufficiently large.  It follows
that, for every required derivative,
\[
 \|\partial^\beta U'_\theta(\cdot,t')\|_{L^2(\R^2)}^2
 \lesssim \sum_{\alpha\in\theta}|c_\alpha|^2.
\]
Time derivatives are included: bounded child frequencies and factors
$a^{-1}$ or $a^{-2}$ from differentiating profiles preserve the same
packet bounds.  Unit-cell Sobolev estimates, bounded overlap of the
enlarged cells, and integration over the $O(1)$ native time layers in
$J$ now give the first assertion of \eqref{eq:src13}.  Summing the
coefficient energies across fine arrays and clusters costs at most $1$.

For the second assertion, let $\mu_\theta=F_{\theta\#}\mu$ in the
coordinates of the relevant cluster.  Covering the inverse image of a
ball shows that its ball parameter is at most $Cmh^3$ for radii at
least $1$.  For every child label, the transverse ratio in
\eqref{eq:src6} is unchanged by \eqref{eq:dyadic-child-coordinates},
and
\begin{equation}\label{eq:dyadic-child-moment}
 (mh^3)a^3/D=mL^3/D.
\end{equation}
Thus $\mu_\theta$ satisfies the same sparsity condition with $D$ after
a fixed inflation of its ball parameter.

Convolve $\mu_\theta$ with $K_P(z')=(1+|z'|)^{-P}$ and restrict the
resulting measure to $|t'|\leq C'a^2$, with a fixed enlargement sufficient
for the cells and derivatives in use.  Call the resulting measure
$\widetilde\mu_\theta$.  Its ball parameter remains $Cmh^3$.  Its
smoothed tube moments are still bounded as in
\eqref{eq:dyadic-child-moment}: translating a child tube kernel by $w$
changes it by at most a factor $C(1+|w|)^M$, and
\[
 \int_{\R^3}(1+|w|)^{M-P}\,dw<\infty.
\]
Every relevant unit cell $q$ in a fixed enlargement of $T$ has
\begin{equation}\label{eq:dyadic-convolved-cell}
 \widetilde\mu_\theta(q)\gtrsim mh^3d(T),
\end{equation}
after using the time enlargement; this follows by comparing the kernel
at all points of $q$ with its distance to $T$.

Apply Proposition~\ref{prop:sparse}, including its derivative version,
at packet size $a$ to $U'_\theta$ and
$\widetilde\mu_\theta$.  Unit cells in $T^*$, of which there are
boundedly many, control the suprema defining $s_{\theta,T}$.
Using \eqref{eq:dyadic-convolved-cell} and dividing out
$(mh^3)^{2/3}$ yields
\[
 \sum_T d(T)^{2/3}s_{\theta,T}^2
 \lesssim h^{-2}a^{4/3}D^{-c_s}
       \sum_{\alpha\in\theta}|c_\alpha|^2.
\]
The transformed spatial domain may exceed size $a^2$.  To use the local
form of the sparse estimate, divide it into boxes of side comparable to
$a^2$ and retain only intersecting compact child packets.  Each such
packet meets boundedly many boxes on $|t'|\leq C'a^2$; summing the
squared estimates therefore uses only bounded total coefficient energy.

This application requires a fixed positive-power lower bound on $D$ in
terms of $a$.  By $a\asymp(BD)^{C_a}$ and $a\geq D$, we have
$a^{1/(2C_a+2)}\leq D\leq a$ once $D$ exceeds a constant depending on
$B$.  Use this fixed exponent in Proposition~\ref{prop:sparse}.
For the remaining bounded values of $D$, also $a$ is bounded in terms
of $B$.  The first assertion summed over $O(a^2)$ time layers, and
$d(T)\lesssim1$, gives the second assertion with a different allowed
constant.  Finally sum over the partitioned coefficient arrays.
\end{proof}

We sort the cells into dyadic density bands
$d(T)\asymp d_\ell$, where $d_\ell=d_{\max}2^{-\ell}$,
$\ell\geq0$, and $d_{\max}$ is a fixed upper bound.  Shifting this
index by a fixed amount has no effect.  Only $O(\log h)$ bands need
occur.  Indeed if the total measure of the domain is already less than
$mL^4D^{-\sigma}g^{-p}$ there is nothing to prove.  Otherwise, for a
relevant cell the transformed domain diameter is $O(L^2/h)$, and the
strictly positive weight in \eqref{eq:dyadic-density} gives
\begin{equation}\label{eq:dyadic-density-lower}
 d(T)\gtrsim (mh^3)^{-1}\mu(\text{domain})
                   (1+CL^2/h)^{-P}
 \gtrsim a^{4-2P}h^{1-P}D^{-\sigma}g^{-p}.
\end{equation}
Here $L=ah$, and fixed powers of $a,D,g$ are bounded by fixed powers of
$X$, which in turn are bounded by fixed powers of $h$.  Thus
$d(T)\geq c_B h^{-C}$ with a fixed $C$.  This band count will only be
used with a power gain in $h$; it will not remain in the final bound.

\subsection{Two bands and the bilinear frame estimate}

For every coarse interval satisfying
\eqref{eq:dyadic-significant-intervals}, summably split its threshold
over the density bands.  At least one band field then has modulus at
least
\begin{equation}\label{eq:dyadic-band-level}
 A_\ell=c\delta_0(1+\ell)^{-2}\lambda.
\end{equation}
The constant $c>0$ is fixed sufficiently small that the sum of these
thresholds is less than the significant-interval threshold.  Choose one
such band for each significant interval. Fix a sufficiently large exponent
$C_d$, to be chosen in the tiny-band estimate below. If one selected band has
$d<X^{-C_d}$, pair its interval with a separated significant interval.
Otherwise there are only $O(C_d\log X)$ selected bands.  By increasing
$B$, the number $c\delta_0^{-\eta}$ in
\eqref{eq:dyadic-significant-intervals} exceeds a sufficiently large
multiple of this band count.  Pigeonholing gives two separated coarse
intervals with the same band.  Here ``separated'' means distance at least
a fixed multiple of $\delta_0$; excluding a bounded number of neighboring
intervals suffices.  The list of significant intervals is also large
enough to supply the partner in the tiny-band case.

Partition time into intervals $J$ of length $h^2$.  For each cluster,
$J$, separated pair $(\tau_1,\tau_2)$, and selected bands
$(\ell_1,\ell_2)$, form the measurable subset where both band fields
exceed the corresponding levels in \eqref{eq:dyadic-band-level}.
Restrict to the high points under consideration if desired.  Denote its
mass by $M_0$, and put
\begin{equation}\label{eq:dyadic-class-energy}
 E_i=\sum_{\substack{\theta\subset\tau_i,\ T\in\mathcal T(J)\\
                         d(T)\asymp d_i}}
                 s_{\theta,T}^2,
 \qquad d_i=d_{\ell_i},\qquad i=1,2.
\end{equation}
The witness subsets need not be disjoint; we use an upper bound on their
sum of masses.  There are at most $X^{C_6}(\log h)^2$ relevant tuples:
the cluster list and interval pairs have polynomial size in $X$, and
the time partition contributes $O(a^2)$ intervals.

Tuples with
\begin{equation}\label{eq:dyadic-small-mass}
 M_0\leq mh^{4-1/12}
\end{equation}
have total mass at most
$X^{C_6}(\log h)^2mh^{4-1/12}$.  Relative to the target this is at most
\[
 C X^{C_6+\sigma+C_gp_{\max}}
          a^{-4}h^{-1/12}(\log h)^2,
\]
which is bounded, and can be made small at large scales, by the choice
of $\kappa$.  This discard counts tuples before normal spatial slabs
are introduced, so it incurs no number of slabs.

Fix a remaining tuple.  Define its two full square envelopes by
\[
 S_i(z)^2=
 \sum_{\substack{\theta\subset\tau_i,\ T\in\mathcal T(J)\\
                         d(T)\asymp d_i}}
 s_{\theta,T}^2(1+\dist(F_\theta(z),T))^{-P}.
\]
By \eqref{eq:dyadic-density},
$\int S_i^2\,d\mu\lesssim mh^3d_iE_i$.  Markov's inequality retains
at least half the witness mass with both bounds
\begin{equation}\label{eq:dyadic-markov}
 S_i\leq B_i,
 \qquad B_i=C\left(\frac{mh^3d_iE_i}{M_0}\right)^{1/2}.
\end{equation}
The constant $C$ here is large and absolute.  The first estimate in
\eqref{eq:src13}, $d_i\lesssim1$, and the reverse of
\eqref{eq:dyadic-small-mass} give
\begin{equation}\label{eq:dyadic-ratios}
 B_i\lesssim h^{-3/2+1/24},\qquad
 \frac{A_{\ell_i}}{B_i}
 \gtrsim h^{1/8}\delta_0g a^{-4/3}(1+\ell_i)^{-2}.
\end{equation}
Since $g\geq D^{-1}$ and $\ell_i=O(\log h)$, the right side of the
ratio estimate is bounded below by $h^{1/8}$ divided by a fixed power
of $X$ and by $C_B(\log h)^2$.  Our scale hierarchy makes it larger
than the fixed threshold-ratio constant in Theorem~\ref{thm:frame}.
The same hierarchy ensures $\delta_0\geq h^{-1/2}$.  Zero class energy
would give an identically zero class field, so it cannot occur in a
tuple with positive witness mass.

\begin{lemma}[Frame estimate for one tuple]\label{lem:dyadic-tuple}
For every tuple remaining after \eqref{eq:dyadic-small-mass},
\begin{equation}\label{eq:dyadic-frame-substitution}
 M_0\lesssim \delta_0^{-C}mh
 \frac{B_1^2B_2^2}{A_{\ell_1}^2A_{\ell_2}^2}
 \left(\frac{h^3E_1}{A_{\ell_1}^2}
            +\frac{h^3E_2}{A_{\ell_2}^2}\right).
\end{equation}
In particular,
\begin{equation}\label{eq:dyadic-cubic-mass}
 \lambda^6M_0^3\lesssim
 \delta_0^{-C}(1+\ell_1+\ell_2)^C
 m^3h^{10}d_1d_2E_1E_2(E_1+E_2).
\end{equation}
\end{lemma}
\begin{proof}
Partition the normal coordinate into slabs of width $h$.  In a slab
centered at $y_0$, only boundedly many normal cell indices $T_y$ can
contribute, and on $J$ only boundedly many time indices $T_t$ can
contribute.  Write $t=t_0+s$ with $t_0$ an endpoint of $J$.
For a composite cell in \eqref{eq:dyadic-composite-packet}, the
tangential carrier has frequency $v_\theta$, and its intercept in the
time coordinate $s$ is
\begin{equation}\label{eq:dyadic-recentered-intercept}
 b_{\theta,T}=hT_x+2v_\theta t_0.
\end{equation}
The fine centers are spaced at scale $h^{-1}$.  At each such center,
$T_x$ runs over an integer lattice translated by $2v_\theta t_0/h$;
the bounded number of normal and time indices thus gives the required
one-dimensional phase-space counts uniformly in $t_0$.  The packet
profiles are compact and smooth on the $h,h^2$ scales, by the definition
of $s_{\theta,T}$.

Expand these profiles, with a fixed smooth cutoff on an enlargement of
the slab, in the variable $(y-y_0)/h$.  The mode packets can be written
with scalar amplitudes at most
\[
 C s_{\theta,T}\langle n\rangle^{-N_1},
 \qquad \langle n\rangle=(1+n^2)^{1/2},
\]
and uniformly bounded required tangential profile derivatives, for a
fixed large $N_1$.  This follows by integrating the Fourier coefficients
by parts, using enough derivatives in the definition of $s$.  The
external Fourier phases have modulus one at every witness.

For the cells belonging to this slab and time bin, the omitted normal
and native time displacements in $F_\theta(z)-T$ are bounded at a
witness.  Hence the one-dimensional envelope weight in
Theorem~\ref{thm:frame}, after \eqref{eq:dyadic-recentered-intercept},
is at most a fixed multiple of the full weight defining $S_i$ there.
On the set retained in \eqref{eq:dyadic-markov}, the $n$th mode
envelope is consequently bounded by
$CB_i\langle n\rangle^{-N_1}$.  Its energy parameter $H_i^*$ in
\eqref{eq:src8} is at most
$Ch^3E_{i,\mathrm{slab}}\langle n\rangle^{-2N_1}$, where
$E_{i,\mathrm{slab}}$ is the part of \eqref{eq:dyadic-class-energy}
near this slab.

Whenever a band field has modulus at least $A_{\ell_i}$, some Fourier
mode has modulus at least $cA_{\ell_i}\langle n\rangle^{-2}$.
For each pair of selected modes apply Theorem~\ref{thm:frame} at
separation $\delta_0$.  The threshold-to-envelope ratios only improve
with $|n|$ when $N_1>2$, so \eqref{eq:dyadic-ratios} verifies its ratio
hypothesis uniformly in the modes.  The factors in the frame estimate
are summable over the two mode indices: for example its first energy
term has powers
$\langle n_1\rangle^{-4N_1+8}
 \langle n_2\rangle^{-2N_1+4}$.
Take $N_1>3$, with sufficient additional regularity for all derivatives.

The projection to $(x,t)$ of the measure restricted to a slab has mass
at most $Cmh$ on any unit square: its inverse image has dimensions
$O(1),O(h),O(1)$ and is covered by $O(h)$ unit balls.  This is an upper
bound on the full projected measure, regardless of which normal
coordinate supplied a witness.  Multiplying the unit-square count from
\eqref{eq:src8} by $Cmh$, and summing the modes, bounds the retained
witness mass in this slab by
\[
 C\delta_0^{-C}mh
 \frac{B_1^2B_2^2}{A_{\ell_1}^2A_{\ell_2}^2}
 \left(\frac{h^3E_{1,\mathrm{slab}}}{A_{\ell_1}^2}
       +\frac{h^3E_{2,\mathrm{slab}}}{A_{\ell_2}^2}\right).
\]
Each composite cell belongs to boundedly many enlarged slabs, and
therefore $\sum_{\mathrm{slab}}E_{i,\mathrm{slab}}\lesssim E_i$.
Summing costs no number of slabs and proves
\eqref{eq:dyadic-frame-substitution}, since at least half of $M_0$
was retained by Markov.

Finally substitute \eqref{eq:dyadic-markov} and
\eqref{eq:dyadic-band-level}.  The product of the two squared envelope
bounds supplies $m^2h^6d_1d_2E_1E_2/M_0^2$; the projected square mass
and frame energy supply $mh\cdot h^3$.  Thus the total factor is
$m^3h^{10}/M_0^2$.  The six threshold powers contribute
$\lambda^{-6}$ times a fixed power of $\delta_0^{-1}$ and of
$1+\ell_1+\ell_2$.  Rearranging proves
\eqref{eq:dyadic-cubic-mass}.
\end{proof}

\subsection{Summing the bands and closing the induction}

First suppose at least one selected band is tiny.  The unweighted
energy bound $E_i\lesssim h^{-2}$ in \eqref{eq:src13} and the cube
root of \eqref{eq:dyadic-cubic-mass} give
\begin{equation}\label{eq:dyadic-tiny-band}
 \lambda^2M_0\lesssim
 \delta_0^{-C}(1+\ell_1+\ell_2)^C
 mh^{4/3}(d_1d_2)^{1/3}.
\end{equation}
For dyadic bands, geometric decay absorbs the polynomial index factors:
\[
 \sum_{\min(d_1,d_2)<X^{-C_d}}
 (1+\ell_1+\ell_2)^C(d_1d_2)^{1/3}
 \lesssim X^{-C_d/6}.
\]
The constants here depend on the fixed exponent $C$.  After summing
over clusters, time bins and coarse interval pairs, all other costs in
\eqref{eq:dyadic-tiny-band} are at most a fixed power of $X$, uniformly
for $u\leq1$.  Choose $C_d$ large enough to absorb this power, together
with $D^\sigma g^{p-2}$ for
\eqref{eq:dyadic-middle-levels}; using $h^{4/3}\leq L^{4/3}$, this
gives
\[
 \lambda^2\sum_{\rm tiny}M_0
 \lesssim mL^{4/3}D^{-\sigma}g^{-(p-2)}.
\]
This is precisely the target multiplied by $\lambda^2$, with a constant
independent of induction.

For the other tuples the two density bands coincide, say $d_1\asymp
d_2\asymp d\geq X^{-C_d}$.  Then $\ell_i=O(C_d\log X)$, and
\[
 \big(E_1E_2(E_1+E_2)\big)^{1/3}\leq C(E_1+E_2).
\]
The cube root of \eqref{eq:dyadic-cubic-mass} gives
\begin{equation}\label{eq:dyadic-common-band}
 \lambda^2M_0\lesssim
 \delta_0^{-C}(\log X)^C
 mh^{10/3}d^{2/3}(E_1+E_2).
\end{equation}
We now sum energies rather than the number of tuples.  A composite
cell appears only in its own cluster, in boundedly many $J$, in its
specified density band and coarse interval, and for at most
$O(\delta_0^{-1})$ choices of the other coarse interval at that band.
There is consequently only an additional fixed power of
$\delta_0^{-1}$ in the multiplicity.  In particular, no factor $a^2$
from time bins, no number of clusters, and no number of density bands
is paid in this sum.  The weighted estimate in \eqref{eq:src13}
therefore gives
\begin{align}
 \lambda^2\sum_{\rm common}M_0
 &\lesssim_B
 X^{C_fu}(\log X)^{C_f}mh^{10/3}
       \big(h^{-2}a^{4/3}D^{-c_s}\big)\notag\\
 &=C_BX^{C_fu}(\log X)^{C_f}
       mL^{4/3}D^{-c_s}.
 \label{eq:dyadic-final-energy}
\end{align}
The fixed exponent $C_f$ includes the just described partner
multiplicity; it is the exponent used in choosing $u$.  The transforms
vary with $\theta$ and with the cluster, but their coefficient arrays
are disjoint, exactly as required in the proof of
\eqref{eq:src13}.

The target after multiplication by $\lambda^2$ is
$mL^{4/3}D^{-\sigma}g^{-(p-2)}$.  Since $g\leq X^{C_g}$, the ratio
of \eqref{eq:dyadic-final-energy} to this quantity is at most
\[
 C_BX^{C_fu+C_g(p-2)}(\log X)^{C_f}D^{\sigma-c_s}.
\]
It is bounded by a fixed constant depending on $B$ by the strict second
inequality in \eqref{eq:dyadic-exponent-conditions}.  This proves the
required estimate for the common-band tuples with no growing loss in
$L$ when $D$ and $g$ are fixed.

\paragraph{Compatible order of parameter choices.}
We now verify that the conditions used in the proof can be imposed together.
The values of $p_{\max}$ and $\eta$ were fixed in
\eqref{eq:dyadic-gamma}. Temporarily restricting $0<u\leq1$ and
$0<\sigma\leq1$ bounds all geometric losses by fixed powers of $X$.
Choose $C_g$ large enough for \eqref{eq:dyadic-large-level}, using only
these powers and the positive gap $p_0-p_{\max}$. Choose $C_a$ larger
still for \eqref{eq:dyadic-plate-discard}, and then choose $C_d$ for the
tiny-band sum. These exponents are independent of how small $u$ will be.

The equal-band estimate \eqref{eq:dyadic-final-energy} has loss
$X^{C_fu}(\log X)^{C_f}$, where the fixed $C_f$ accounts for the frame
separation loss and the coarse-interval multiplicity. Choose $u>0$ so
small that $C_fu<c_s/4$. Next choose $\kappa>0$ small in terms of all
the fixed powers just chosen. For any fixed required $K$, the relations
$X\leq BL^\kappa$ and $h\asymp L/X^{C_a}$ give $h\geq X^K$ at
sufficiently large $L$ if $\kappa(C_a+K)<1$. This supplies the geometric
scale conditions, $\delta_0\geq h^{-1/2}$, the frame threshold ratios,
and the small-mass discard.

Choose $p>2$ sufficiently close to $2$, with $p\leq p_{\max}$, and
then choose $\sigma>0$ sufficiently small to satisfy all three strict
inequalities in \eqref{eq:dyadic-exponent-conditions}. Finally choose
$B$ large enough to contract the recursive terms and to supply the angular
pigeonhole argument. Constants depending on $B$ are allowed in the
nonrecursive estimates. Enlarge the fixed base scale $L_0$ to cover the
remaining comparisons involving $B$. Every fixed constant and exponent is
chosen before the induction constant and the variable scales $L,D$.

\begin{proof}[Completion of the proof of Theorem~\ref{thm:distribution}]
The sparse range and bounded scales were handled directly.  In the
remaining range, the prefix exceptions, all dyadic-cap exceptions, and
the coarse cluster exceptions cost less than half the desired bound
with constant $C_{\rm ind}$, by
\eqref{eq:dyadic-contraction}.  Every other original high point either
is handled by the large-level strip estimate or total mass, lies in a
discarded plate intersection, belongs to a small-mass tuple, or is
covered by the tiny-band or common-band bounds.  All these latter costs
are bounded by a fixed multiple of $mL^4D^{-\sigma}g^{-p}$ independently
of $C_{\rm ind}$.  Choose $C_{\rm ind}$ once to be twice their total
constant and large enough for the initial cases.  The two inductions
close.  Homogeneity supplies the factor $\|c\|_{\ell^2}^p$ for arbitrary
subarrays and completes the theorem.
\end{proof}
